\section{Background and Setup}
\label{sec:background}

\paragraph{RAG and faithfulness.}
A RAG pipeline retrieves $k$ passages
$\{p_1, \ldots, p_k\}$ for a query $q$, concatenates them as
context, and asks a generator $\mathcal{G}$ to produce an answer
$a = \mathcal{G}(q, p_1, \ldots, p_k)$. Faithfulness is the question
of whether the claims in $a$ are entailed by
$\{p_1, \ldots, p_k\}$. Faithfulness is typically operationalized
by an entailment scorer $\phi: (\text{premise},
\text{hypothesis}) \to [0, 1]$ run sentence-by-sentence and
averaged. The dominant choice in recent RAG papers is
\texttt{cross-encoder/nli-deberta-v3-base}
\citep{he2021debertav3}.

\paragraph{The refinement paradox.}
A line of recent work \citep{xu2023recomp} examines what happens
when the retrieved set is post-processed to retain only
high-relevance passages. The intuition is that filtering out
irrelevant retrievals should improve faithfulness. \emph{Hybrid
Context-Preserving Chunking} (HCPC) and its
coherence-aware variant HCPC-v2, introduced as a controlled
intervention probe, have been reported to expose a counterintuitive
\emph{refinement paradox}: aggressive refinement that maximizes
per-passage similarity reduces answer faithfulness, while a
selective coherence-aware refinement (HCPC-v2) recovers it. Earlier
results at $n{=}30$ per condition reported a SQuAD/Mistral-7B
paradox magnitude of $0.069 \pm 0.004$ across three seeds, with a
signal-to-standard-deviation ratio of $\sim 17$.

\paragraph{Context coherence as a candidate mediator.}
The proposed mechanism is the \emph{Context Coherence Score} (CCS),
a generator-free retrieval-time diagnostic defined as
\(\mathrm{CCS}(\{p_1, \ldots, p_k\}) = \overline{\mathrm{sim}}_{ij}
- \sigma_{\mathrm{sim}_{ij}}\), where the average and standard
deviation are taken over the off-diagonal pairwise embedding
similarities of the retrieved passages.
The hypothesized causal claim is that CCS, not per-passage
relevance, mediates faithfulness: at fixed mean query-passage
similarity, lower CCS would produce lower faithfulness.

\paragraph{Why this needs auditing.}
Three properties of the published evidence make a pre-registered
audit appropriate. First, the headline statistic comes from
$n{=}30$ queries per condition; small samples are notorious for
inflating effect sizes. Second, faithfulness is reported under a
single scorer (DeBERTa-v3 NLI); we have no published cross-scorer
agreement statistics on RAG outputs at scale. Third, the threshold
$\tau$ that gates HCPC's refinement decision is tuned on a 50-query
held-out subset of SQuAD --- the same dataset the paradox is
reported on. The latter two are common practices across the broader
RAG literature, not unique to this work. We use the refinement
paradox as a stress test of those practices.

\paragraph{Pre-registration.}
For each experiment we fix sample sizes, seeds, scoring code, and
statistical tests \emph{before} running any generation, in
machine-checkable logs at
\texttt{experiments/fix\_NN\_log.md}. Effect-size thresholds and
honest-reporting commitments (what we will say if a hypothesis is
not supported) are also fixed in advance. We use paired Wilcoxon
signed-rank for paired contrasts, 10{,}000-resample bootstrap CIs
for continuous summary statistics, and Wilson score CIs for binary
proportions. Cohen's $d_z$ is reported alongside every paired
contrast, regardless of significance.
